% !TeX root = ../main.tex

%\chapter{Summary}

Two physics analyses with a $\ell^+\ell^- + E_\text{T}^\text{miss}$ final state are presented in this dissertation. 
Both analyses use the full Run~2 dataset corresponding to an integrated luminosity of 139~fb$^{-1}$, collected by the ATLAS experiment in proton--proton collisions at a center-of-mass energy of $\sqrt{s}=13$~TeV during the LHC Run~2 program.

The first analysis presents a precision measurement of the total and differential cross-sections of the $pp \to ZZ \to \ell^+\ell^-\nu\bar{\nu}$ process. 
The measured cross-sections are found to be in good agreement with Standard Model (SM) predictions. 
The measurement achieves a precision of 4.6\%, representing more than a 50\% improvement compared to the previous precision of 7\%. 
These results lead to significantly more stringent constraints on anomalous triple- and quartic-gauge-boson self-interactions, providing a precise test of the electroweak sector of the Standard Model.

The second analysis focuses on a search for DM particles in the mono-$Z$ channel. 
No significant deviation from the SM prediction is observed, and exclusion limits are derived within several theoretical frameworks. 
The search sensitivity is significantly improved compared to previous results in the same channel, benefiting from the larger dataset and improved analysis techniques.

In addition to the physics analyses, this dissertation also reports on the construction and testing of sMDT chambers for the ATLAS detector upgrade for the HL-LHC. 
The performance of the newly constructed sMDT chambers is evaluated using cosmic-ray data, demonstrating excellent tracking resolution and high detection efficiency, consistent with the requirements for future ATLAS operation.

Looking forward, these physics topics will continue to benefit from the much larger datasets collected in Run~3 and at the HL-LHC. 
The more than 20-fold increase in integrated luminosity will enable precision measurements and searches at the LHC to provide even more stringent tests of the Standard Model and enhance the discovery potential for physics beyond it. 
This dissertation establishes a solid foundation for future studies in these topics.

The sMDT chambers will be installed in the ATLAS detector in 2028--2029. 
They will significantly enhance the ATLAS muon trigger capabilities at Level-1 and allow operation under substantially higher background rates associated with the increased LHC luminosity, ensuring efficient and high-quality muon data collection and reconstruction for physics analyses.